\documentclass[11pt,a4paper]{article}

\usepackage[margin=1in]{geometry}
\usepackage{setspace}
\setstretch{1.1}
\usepackage[T1]{fontenc}
\usepackage{lmodern}
\usepackage[hidelinks]{hyperref}
\usepackage{parskip}

\pagestyle{empty}

\newcommand{\axionbloch}{\texttt{axionbloch}}

\begin{document}
\begin{flushright}
Yuzhe Zhang \\
Helmholtz Institute Mainz \\
Staudingerweg 18 \\
55128 Mainz, Germany \\
\href{mailto:yuhzhang@uni-mainz.de}{yuhzhang@uni-mainz.de} \\
\today
\end{flushright}

Dear Editors of \textit{Universe},

Please consider my manuscript, ``\axionbloch: an Open-Source Python Package for Simulating Axion-Induced Spin Dynamics,'' for publication in \textit{Universe} as part of the Special Issue ``Ultralight Bosonic Dark Matter: Theoretical Developments and Experimental Searches.''

The manuscript presents \axionbloch{}, an open-source Python package for simulating spin dynamics induced by axion and other pseudomagnetic fields. The work is significant because it provides a practical numerical tool for studying axion-induced signals in nuclear magnetic resonance and related precision-measurement settings, where reliable simulation and calibration are essential for interpreting data and designing searches. By combining a physically transparent Bloch-equation framework with implementation details, calibration tests, and example simulations, the paper offers both methodological value and immediate utility for the ultralight dark matter community.

The manuscript fits the scope of \textit{Universe} and this Special Issue because it addresses ultralight bosonic dark matter, axionlike particles, and experimental searches for exotic spin-dependent interactions. I believe it will be of interest to readers working on axion theory, magnetic resonance, and precision measurement techniques.

As the only author of the manuscript, I confirm that neither the manuscript nor any parts of its content are currently under consideration for publication with or published in another journal.
I have approved the manuscript and agree with its submission to \textit{Universe}.

\vspace{1.5em}

\noindent Yours sincerely,

\noindent Yuzhe Zhang

\end{document}
